% ------------------------------------------------------------------------
%  Linear solver comparison (bench_sweep.py sweep A).
%  Generated by scripts/plot_solver_bench.py -- regenerate rather than edit.
%  Include with % ------------------------------------------------------------------------
%  Linear solver comparison (bench_sweep.py sweep A).
%  Generated by scripts/plot_solver_bench.py -- regenerate rather than edit.
%  Include with % ------------------------------------------------------------------------
%  Linear solver comparison (bench_sweep.py sweep A).
%  Generated by scripts/plot_solver_bench.py -- regenerate rather than edit.
%  Include with % ------------------------------------------------------------------------
%  Linear solver comparison (bench_sweep.py sweep A).
%  Generated by scripts/plot_solver_bench.py -- regenerate rather than edit.
%  Include with \input{figures/solver_bench_cluster}
%  Needs \usepackage{pgfplots} in the preamble.
% ------------------------------------------------------------------------
\input{figures/data/solver_bench_cluster_ref}

\pgfplotsset{compat=1.16}

% Keep only the rows whose column #1 equals #2.
\pgfplotsset{
  discard if not/.style 2 args={
    x filter/.code={%
      \edef\tempa{\thisrow{#1}}\edef\tempb{#2}%
      \ifx\tempa\tempb\else\def\pgfmathresult{inf}\fi
    }
  }
}

\definecolor{solverlis}{HTML}{D95F02}
\definecolor{solveramgx}{HTML}{9DC3E0}
\definecolor{solvercfg}{HTML}{2C5F8D}

\begin{figure}[!htb]
\centering
\begin{tikzpicture}
\begin{axis}[
    ybar,
    % The iteration-count label sits above the top whisker by 5pt, which
    % ymax's headroom does not always cover -- clip=false stops that label from
    % being cut off instead of chasing a headroom multiplier that varies with
    % how many digits the label has.
    clip=false,
    width=13.1cm, height=6.4cm,
    bar width=10pt,
    enlarge x limits=0.28,
    ymode=log, log origin y=infty,
    ymin=0.759405, ymax=17.3883,
    ytick={1,2,5,10},
    log ticks with fixed point,
    ylabel={\solverMetricsolverbenchcluster{} time, relative to the fastest},
    ylabel style={font=\footnotesize},
    symbolic x coords={6VYB,1VSZ,H1N1},
    xtick=data,
    xticklabel style={font=\footnotesize, align=center},
    tick label style={font=\scriptsize},
    ymajorgrids=true,
    major grid style={black!12},
    legend style={
      at={(0.5,-0.16)}, anchor=north, draw=none,
      legend columns=3, font=\footnotesize, column sep=0.8em,
    },
    % Run-to-run spread is only a few percent, so the upper whisker is 1-2 mm and
    % the lower half is buried in the bar fill. Caps nearly as wide as the bar
    % are what make it read as an error bar rather than a stray tick.
    % T-caps. rotate=90 is required: pgfplots draws the cap perpendicular to the
    % error direction by itself, but any mark size given here clobbers that, and
    % the cap then renders vertically -- which reads as a longer whisker rather
    % than a wider cap. Setting error mark explicitly does not help; the rotation
    % is what matters.
    error bars/error bar style={black!85, line width=0.8pt},
    error bars/error mark options={
      mark size=2.5pt, rotate=90, line width=0.8pt, black!85,
    },
    every node near coord/.append style={
      font=\tiny, color=black!60, anchor=south, yshift=5pt,
    },
]
  \addplot+[
    ybar, fill=solverlis, draw=solverlis!70!black, line width=0.4pt,
    nodes near coords, point meta=explicit symbolic,
    error bars/.cd, y dir=both, y explicit,
  ] table[x=molecule, y=lis, y error minus=lis_em, y error plus=lis_ep,
          meta=lis_it] {figures/data/solver_bench_cluster.dat};
  \addlegendentry{LIS}
  \addplot+[
    ybar, fill=solveramgx, draw=solveramgx!70!black, line width=0.4pt,
    nodes near coords, point meta=explicit symbolic,
    error bars/.cd, y dir=both, y explicit,
  ] table[x=molecule, y=amgx, y error minus=amgx_em, y error plus=amgx_ep,
          meta=amgx_it] {figures/data/solver_bench_cluster.dat};
  \addlegendentry{AMGX (default)}
  \addplot+[
    ybar, fill=solvercfg, draw=solvercfg!70!black, line width=0.4pt,
    nodes near coords, point meta=explicit symbolic,
    error bars/.cd, y dir=both, y explicit,
  ] table[x=molecule, y=amgxcfg, y error minus=amgxcfg_em, y error plus=amgxcfg_ep,
          meta=amgxcfg_it] {figures/data/solver_bench_cluster.dat};
  \addlegendentry{AMGX (tuned)}

\end{axis}
\end{tikzpicture}

\caption[Linear solver comparison]{\textbf{Linear solver comparison on
    \solverMoleculessolverbenchcluster{}.} Median \solverMetricsolverbenchcluster{} time over \solverRepeatssolverbenchcluster{} runs on a logarithmic axis,
    whiskers spanning fastest to slowest; the figure above each column is the
    iteration count.\solverMachinesolverbenchcluster{}\solverCaveatssolverbenchcluster{}}
\label{fig:solver-bench-cluster}
\end{figure}

%  Needs \usepackage{pgfplots} in the preamble.
% ------------------------------------------------------------------------
% generated by scripts/plot_solver_bench.py -- do not edit
\def\solverMetricsolverbenchcluster{solve stage}
\def\solverRepeatssolverbenchcluster{3}
\def\solverUnitsolverbenchcluster{relative to the fastest}
\def\solverMoleculessolverbenchcluster{6VYB, 1VSZ, H1N1}
\def\solverEnergyMethodsolverbenchcluster{1}
\def\solverMachinesolverbenchcluster{ Measured on an A100 node of the TU Berlin HPC (Table~\ref{tab:test-systems}) with 16 MPI ranks.}
\def\solverCaveatssolverbenchcluster{ H1N1: measured on four H200 GPUs of the TU Berlin HPC (Table~\ref{tab:test-systems}) with 8 MPI ranks; 1 run per configuration, not 3.}


\pgfplotsset{compat=1.16}

% Keep only the rows whose column #1 equals #2.
\pgfplotsset{
  discard if not/.style 2 args={
    x filter/.code={%
      \edef\tempa{\thisrow{#1}}\edef\tempb{#2}%
      \ifx\tempa\tempb\else\def\pgfmathresult{inf}\fi
    }
  }
}

\definecolor{solverlis}{HTML}{D95F02}
\definecolor{solveramgx}{HTML}{9DC3E0}
\definecolor{solvercfg}{HTML}{2C5F8D}

\begin{figure}[!htb]
\centering
\begin{tikzpicture}
\begin{axis}[
    ybar,
    % The iteration-count label sits above the top whisker by 5pt, which
    % ymax's headroom does not always cover -- clip=false stops that label from
    % being cut off instead of chasing a headroom multiplier that varies with
    % how many digits the label has.
    clip=false,
    width=13.1cm, height=6.4cm,
    bar width=10pt,
    enlarge x limits=0.28,
    ymode=log, log origin y=infty,
    ymin=0.759405, ymax=17.3883,
    ytick={1,2,5,10},
    log ticks with fixed point,
    ylabel={\solverMetricsolverbenchcluster{} time, relative to the fastest},
    ylabel style={font=\footnotesize},
    symbolic x coords={6VYB,1VSZ,H1N1},
    xtick=data,
    xticklabel style={font=\footnotesize, align=center},
    tick label style={font=\scriptsize},
    ymajorgrids=true,
    major grid style={black!12},
    legend style={
      at={(0.5,-0.16)}, anchor=north, draw=none,
      legend columns=3, font=\footnotesize, column sep=0.8em,
    },
    % Run-to-run spread is only a few percent, so the upper whisker is 1-2 mm and
    % the lower half is buried in the bar fill. Caps nearly as wide as the bar
    % are what make it read as an error bar rather than a stray tick.
    % T-caps. rotate=90 is required: pgfplots draws the cap perpendicular to the
    % error direction by itself, but any mark size given here clobbers that, and
    % the cap then renders vertically -- which reads as a longer whisker rather
    % than a wider cap. Setting error mark explicitly does not help; the rotation
    % is what matters.
    error bars/error bar style={black!85, line width=0.8pt},
    error bars/error mark options={
      mark size=2.5pt, rotate=90, line width=0.8pt, black!85,
    },
    every node near coord/.append style={
      font=\tiny, color=black!60, anchor=south, yshift=5pt,
    },
]
  \addplot+[
    ybar, fill=solverlis, draw=solverlis!70!black, line width=0.4pt,
    nodes near coords, point meta=explicit symbolic,
    error bars/.cd, y dir=both, y explicit,
  ] table[x=molecule, y=lis, y error minus=lis_em, y error plus=lis_ep,
          meta=lis_it] {figures/data/solver_bench_cluster.dat};
  \addlegendentry{LIS}
  \addplot+[
    ybar, fill=solveramgx, draw=solveramgx!70!black, line width=0.4pt,
    nodes near coords, point meta=explicit symbolic,
    error bars/.cd, y dir=both, y explicit,
  ] table[x=molecule, y=amgx, y error minus=amgx_em, y error plus=amgx_ep,
          meta=amgx_it] {figures/data/solver_bench_cluster.dat};
  \addlegendentry{AMGX (default)}
  \addplot+[
    ybar, fill=solvercfg, draw=solvercfg!70!black, line width=0.4pt,
    nodes near coords, point meta=explicit symbolic,
    error bars/.cd, y dir=both, y explicit,
  ] table[x=molecule, y=amgxcfg, y error minus=amgxcfg_em, y error plus=amgxcfg_ep,
          meta=amgxcfg_it] {figures/data/solver_bench_cluster.dat};
  \addlegendentry{AMGX (tuned)}

\end{axis}
\end{tikzpicture}

\caption[Linear solver comparison]{\textbf{Linear solver comparison on
    \solverMoleculessolverbenchcluster{}.} Median \solverMetricsolverbenchcluster{} time over \solverRepeatssolverbenchcluster{} runs on a logarithmic axis,
    whiskers spanning fastest to slowest; the figure above each column is the
    iteration count.\solverMachinesolverbenchcluster{}\solverCaveatssolverbenchcluster{}}
\label{fig:solver-bench-cluster}
\end{figure}

%  Needs \usepackage{pgfplots} in the preamble.
% ------------------------------------------------------------------------
% generated by scripts/plot_solver_bench.py -- do not edit
\def\solverMetricsolverbenchcluster{solve stage}
\def\solverRepeatssolverbenchcluster{3}
\def\solverUnitsolverbenchcluster{relative to the fastest}
\def\solverMoleculessolverbenchcluster{6VYB, 1VSZ, H1N1}
\def\solverEnergyMethodsolverbenchcluster{1}
\def\solverMachinesolverbenchcluster{ Measured on an A100 node of the TU Berlin HPC (Table~\ref{tab:test-systems}) with 16 MPI ranks.}
\def\solverCaveatssolverbenchcluster{ H1N1: measured on four H200 GPUs of the TU Berlin HPC (Table~\ref{tab:test-systems}) with 8 MPI ranks; 1 run per configuration, not 3.}


\pgfplotsset{compat=1.16}

% Keep only the rows whose column #1 equals #2.
\pgfplotsset{
  discard if not/.style 2 args={
    x filter/.code={%
      \edef\tempa{\thisrow{#1}}\edef\tempb{#2}%
      \ifx\tempa\tempb\else\def\pgfmathresult{inf}\fi
    }
  }
}

\definecolor{solverlis}{HTML}{D95F02}
\definecolor{solveramgx}{HTML}{9DC3E0}
\definecolor{solvercfg}{HTML}{2C5F8D}

\begin{figure}[!htb]
\centering
\begin{tikzpicture}
\begin{axis}[
    ybar,
    % The iteration-count label sits above the top whisker by 5pt, which
    % ymax's headroom does not always cover -- clip=false stops that label from
    % being cut off instead of chasing a headroom multiplier that varies with
    % how many digits the label has.
    clip=false,
    width=13.1cm, height=6.4cm,
    bar width=10pt,
    enlarge x limits=0.28,
    ymode=log, log origin y=infty,
    ymin=0.759405, ymax=17.3883,
    ytick={1,2,5,10},
    log ticks with fixed point,
    ylabel={\solverMetricsolverbenchcluster{} time, relative to the fastest},
    ylabel style={font=\footnotesize},
    symbolic x coords={6VYB,1VSZ,H1N1},
    xtick=data,
    xticklabel style={font=\footnotesize, align=center},
    tick label style={font=\scriptsize},
    ymajorgrids=true,
    major grid style={black!12},
    legend style={
      at={(0.5,-0.16)}, anchor=north, draw=none,
      legend columns=3, font=\footnotesize, column sep=0.8em,
    },
    % Run-to-run spread is only a few percent, so the upper whisker is 1-2 mm and
    % the lower half is buried in the bar fill. Caps nearly as wide as the bar
    % are what make it read as an error bar rather than a stray tick.
    % T-caps. rotate=90 is required: pgfplots draws the cap perpendicular to the
    % error direction by itself, but any mark size given here clobbers that, and
    % the cap then renders vertically -- which reads as a longer whisker rather
    % than a wider cap. Setting error mark explicitly does not help; the rotation
    % is what matters.
    error bars/error bar style={black!85, line width=0.8pt},
    error bars/error mark options={
      mark size=2.5pt, rotate=90, line width=0.8pt, black!85,
    },
    every node near coord/.append style={
      font=\tiny, color=black!60, anchor=south, yshift=5pt,
    },
]
  \addplot+[
    ybar, fill=solverlis, draw=solverlis!70!black, line width=0.4pt,
    nodes near coords, point meta=explicit symbolic,
    error bars/.cd, y dir=both, y explicit,
  ] table[x=molecule, y=lis, y error minus=lis_em, y error plus=lis_ep,
          meta=lis_it] {figures/data/solver_bench_cluster.dat};
  \addlegendentry{LIS}
  \addplot+[
    ybar, fill=solveramgx, draw=solveramgx!70!black, line width=0.4pt,
    nodes near coords, point meta=explicit symbolic,
    error bars/.cd, y dir=both, y explicit,
  ] table[x=molecule, y=amgx, y error minus=amgx_em, y error plus=amgx_ep,
          meta=amgx_it] {figures/data/solver_bench_cluster.dat};
  \addlegendentry{AMGX (default)}
  \addplot+[
    ybar, fill=solvercfg, draw=solvercfg!70!black, line width=0.4pt,
    nodes near coords, point meta=explicit symbolic,
    error bars/.cd, y dir=both, y explicit,
  ] table[x=molecule, y=amgxcfg, y error minus=amgxcfg_em, y error plus=amgxcfg_ep,
          meta=amgxcfg_it] {figures/data/solver_bench_cluster.dat};
  \addlegendentry{AMGX (tuned)}

\end{axis}
\end{tikzpicture}

\caption[Linear solver comparison]{\textbf{Linear solver comparison on
    \solverMoleculessolverbenchcluster{}.} Median \solverMetricsolverbenchcluster{} time over \solverRepeatssolverbenchcluster{} runs on a logarithmic axis,
    whiskers spanning fastest to slowest; the figure above each column is the
    iteration count.\solverMachinesolverbenchcluster{}\solverCaveatssolverbenchcluster{}}
\label{fig:solver-bench-cluster}
\end{figure}

%  Needs \usepackage{pgfplots} in the preamble.
% ------------------------------------------------------------------------
% generated by scripts/plot_solver_bench.py -- do not edit
\def\solverMetricsolverbenchcluster{solve stage}
\def\solverRepeatssolverbenchcluster{3}
\def\solverUnitsolverbenchcluster{relative to the fastest}
\def\solverMoleculessolverbenchcluster{6VYB, 1VSZ, H1N1}
\def\solverEnergyMethodsolverbenchcluster{1}
\def\solverMachinesolverbenchcluster{ Measured on an A100 node of the TU Berlin HPC (Table~\ref{tab:test-systems}) with 16 MPI ranks.}
\def\solverCaveatssolverbenchcluster{ H1N1: measured on four H200 GPUs of the TU Berlin HPC (Table~\ref{tab:test-systems}) with 8 MPI ranks; 1 run per configuration, not 3.}


\pgfplotsset{compat=1.16}

% Keep only the rows whose column #1 equals #2.
\pgfplotsset{
  discard if not/.style 2 args={
    x filter/.code={%
      \edef\tempa{\thisrow{#1}}\edef\tempb{#2}%
      \ifx\tempa\tempb\else\def\pgfmathresult{inf}\fi
    }
  }
}

\definecolor{solverlis}{HTML}{D95F02}
\definecolor{solveramgx}{HTML}{9DC3E0}
\definecolor{solvercfg}{HTML}{2C5F8D}

\begin{figure}[!htb]
\centering
\begin{tikzpicture}
\begin{axis}[
    ybar,
    % The iteration-count label sits above the top whisker by 5pt, which
    % ymax's headroom does not always cover -- clip=false stops that label from
    % being cut off instead of chasing a headroom multiplier that varies with
    % how many digits the label has.
    clip=false,
    width=13.1cm, height=6.4cm,
    bar width=10pt,
    enlarge x limits=0.28,
    ymode=log, log origin y=infty,
    ymin=0.759405, ymax=17.3883,
    ytick={1,2,5,10},
    log ticks with fixed point,
    ylabel={\solverMetricsolverbenchcluster{} time, relative to the fastest},
    ylabel style={font=\footnotesize},
    symbolic x coords={6VYB,1VSZ,H1N1},
    xtick=data,
    xticklabel style={font=\footnotesize, align=center},
    tick label style={font=\scriptsize},
    ymajorgrids=true,
    major grid style={black!12},
    legend style={
      at={(0.5,-0.16)}, anchor=north, draw=none,
      legend columns=3, font=\footnotesize, column sep=0.8em,
    },
    % Run-to-run spread is only a few percent, so the upper whisker is 1-2 mm and
    % the lower half is buried in the bar fill. Caps nearly as wide as the bar
    % are what make it read as an error bar rather than a stray tick.
    % T-caps. rotate=90 is required: pgfplots draws the cap perpendicular to the
    % error direction by itself, but any mark size given here clobbers that, and
    % the cap then renders vertically -- which reads as a longer whisker rather
    % than a wider cap. Setting error mark explicitly does not help; the rotation
    % is what matters.
    error bars/error bar style={black!85, line width=0.8pt},
    error bars/error mark options={
      mark size=2.5pt, rotate=90, line width=0.8pt, black!85,
    },
    every node near coord/.append style={
      font=\tiny, color=black!60, anchor=south, yshift=5pt,
    },
]
  \addplot+[
    ybar, fill=solverlis, draw=solverlis!70!black, line width=0.4pt,
    nodes near coords, point meta=explicit symbolic,
    error bars/.cd, y dir=both, y explicit,
  ] table[x=molecule, y=lis, y error minus=lis_em, y error plus=lis_ep,
          meta=lis_it] {figures/data/solver_bench_cluster.dat};
  \addlegendentry{LIS}
  \addplot+[
    ybar, fill=solveramgx, draw=solveramgx!70!black, line width=0.4pt,
    nodes near coords, point meta=explicit symbolic,
    error bars/.cd, y dir=both, y explicit,
  ] table[x=molecule, y=amgx, y error minus=amgx_em, y error plus=amgx_ep,
          meta=amgx_it] {figures/data/solver_bench_cluster.dat};
  \addlegendentry{AMGX (default)}
  \addplot+[
    ybar, fill=solvercfg, draw=solvercfg!70!black, line width=0.4pt,
    nodes near coords, point meta=explicit symbolic,
    error bars/.cd, y dir=both, y explicit,
  ] table[x=molecule, y=amgxcfg, y error minus=amgxcfg_em, y error plus=amgxcfg_ep,
          meta=amgxcfg_it] {figures/data/solver_bench_cluster.dat};
  \addlegendentry{AMGX (tuned)}

\end{axis}
\end{tikzpicture}

\caption[Linear solver comparison]{\textbf{Linear solver comparison on
    \solverMoleculessolverbenchcluster{}.} Median \solverMetricsolverbenchcluster{} time over \solverRepeatssolverbenchcluster{} runs on a logarithmic axis,
    whiskers spanning fastest to slowest; the figure above each column is the
    iteration count.\solverMachinesolverbenchcluster{}\solverCaveatssolverbenchcluster{}}
\label{fig:solver-bench-cluster}
\end{figure}
